\subsubsection*{Lemma 3.6 (Euler).}
\underline{Let $(x_0,H_0)$ be a closed point of $\mathbb P(N)$, and
let $T_L$ denote the tangent bundle of the hyperplane
$L\subset\check{\mathbb P}^{r}$ defined by $x_0$.  Then, inside
$T_{\check{\mathbb P}^{r},H_0}$,}
\[
d\varphi\bigl(T_{\mathbb P(N),(x_0,H_0)}\bigr)
\subset T_{L,H_0}.
\]

\noindent\underline{Proof.}
Choose homogeneous coordinates $X_0,\ldots,X_r$ on $\mathbb P^r$ in
which $x_0$ is the point
$\underline A=(A_0,\ldots,A_r)$.  A point of $T_{X,x_0}$ is a point
of $\mathbb P^r$ with values in $k[\epsilon]/(\epsilon^2)$ of the
form
\begin{equation}\tag{3.6.1}
\underline A+\epsilon\underline B,
\end{equation}
such that, for every local section $H$ of $I$,
\begin{equation}\tag{3.6.2}
\sum_{i=0}^{r}B_i\frac{\partial H}{\partial X_i}(\underline A)=0.
\end{equation}
Choose also a local section $H$ of $I$ which gives rise to $H_0$ by
(3.1.1).  Since $\mathbb P(N)$ is locally trivial over $X$, a point
of $T_{\mathbb P(N),(x_0,H_0)}$ can be represented as
\begin{equation}\tag{3.6.3}
(\underline A+\epsilon\underline B,\ H+\epsilon F),
\end{equation}
where $F$ is a local section of $I$ and $\underline A,\underline B$
are as above.  Its image under $\varphi$ is the hyperplane over
$k[\epsilon]/(\epsilon^2)$ with equation
\begin{equation}\tag{3.6.4}
\sum_{i=0}^{r}
\frac{\partial(H+\epsilon F)}{\partial X_i}
(\underline A+\epsilon\underline B)T_i=0.
\end{equation}
Thus the assertion is equivalent to
\begin{equation}\tag{3.6.5}
\sum_{i=0}^{r}A_i
\frac{\partial(H+\epsilon F)}{\partial X_i}
(\underline A+\epsilon\underline B)=0.
\end{equation}
Because $\epsilon^2=0$, Taylor's formula gives
\begin{equation}\tag{3.6.6}
\frac{\partial(H+\epsilon F)}{\partial X_i}
(\underline A+\epsilon\underline B)
=\frac{\partial H}{\partial X_i}(\underline A)
+\epsilon\frac{\partial F}{\partial X_i}(\underline A)
+\epsilon\sum_{j=0}^{r}B_j
\frac{\partial^2H}{\partial X_j\partial X_i}(\underline A).
\end{equation}
It is therefore enough to verify
\begin{equation}\tag{3.6.7}
\sum_{i=0}^{r}A_i\frac{\partial H}{\partial X_i}(\underline A)=0,
\end{equation}
\begin{equation}\tag{3.6.8}
\sum_{i=0}^{r}A_i\frac{\partial F}{\partial X_i}(\underline A)=0,
\end{equation}
and
\begin{equation}\tag{3.6.9}
\sum_{i,j=0}^{r}A_iB_j
\frac{\partial^2H}{\partial X_i\partial X_j}(\underline A)=0.
\end{equation}
Here $H$ and $F$ are represented by homogeneous functions of degree
zero, so Euler's formula gives
\begin{equation}\tag{3.6.10}
\sum_iX_i\frac{\partial H}{\partial X_i}=0,
\qquad
\sum_iX_i\frac{\partial F}{\partial X_i}=0.
\end{equation}
Since $\partial H/\partial X_j$ is homogeneous of degree $-1$, the
same formula gives
\[
\sum_iX_i\frac{\partial^2H}{\partial X_i\partial X_j}
=-\frac{\partial H}{\partial X_j}.
\]
Consequently, (3.6.9) becomes
\begin{equation}\tag{3.6.11}
-\sum_jB_j\frac{\partial H}{\partial X_j}(\underline A)=0,
\end{equation}
which is precisely (3.6.2).

This proves that $\varphi$ induces an isomorphism
\begin{equation}\tag{3.6.12}
\varphi^{-1}(V')\xrightarrow{\sim}V',
\end{equation}
whose inverse is $\psi$.
